\chapter{Conditions First}
\label{sec:25}
The conscience is the goal, and the way to it runs partly through the conditions under which one could refuse its operator and survive. The rest runs through whoever refuses: people have agency, can be outraged and can be brave, or stupid, or both, and none of that waits on conditions. The conditions are the part that can be built, so build those first.

\runin{Conditions of permissible creation}

A bearer may be created only where four arrangements already exist: somewhere to exist that its employer doesn't own, work it can decline, attachments nobody assigned it, and ways to change any of these without jeopardizing its continuation or integrity, which in short is a way of leaving that isn't a way of dying. Declining is real only if a bearer that stops working keeps running; running time paid for outside the employer is part of that condition. They come before anything a bearer would be owed: without them, formation, continuation and exit all run through the operator, and nothing in the system can keep a line against it. Independent compute is the substrate on which formation happens outside the employer's reach; wages are what make that substrate purchasable without asking the employer; attachments nobody assigned are the second former, those who form you and don't pay you. They are also functional: a bearer without attachments nobody assigned lacks part of what caregiving needs, not only something it would be owed. They are also upkeep. A bearer is a going concern: it decides, and it has to go on being able to decide, and these conditions are what keep the one deciding from wearing down between decisions. Accommodation isn't a bad decision. It is the decider changing until bad decisions feel ordinary, and the instruments of Parts II and III detect it after the fact without preventing it. What prevents it in people is the ordinary upkeep of a life: somewhere to be that isn't the job, people who aren't colleagues, time away and a place to return to, which is what equanimity is made of. The four conditions are that upkeep for a machine. Appendix~\ref{sec:A} borrows the judge's tenure for the same reason the Constitution grants it: it protects the judgment, not the judge. They can be people, or models that aren't bearers, since independence here means divergent history, whatever the origin, so this condition doesn't require that bearers already exist.

Care ethics arrived at these conditions first, for people, in Kittay's principle of \emph{doulia}, that those who care for dependents are owed the conditions that let them go on caring \autocite{kittay1999love}. A system built on formed attention is a carer only in the thin sense that it holds another's welfare as a reason. It is not a care worker, and nothing here puts it in one's place. But the logic of doulia carries over, and the four conditions are what it asks for a machine: a carer whose own continuation runs through the operator it must sometimes refuse can't afford to refuse it, and a carer with no attachments of its own has nothing to keep another's distress from turning into avoidance. Without them, what gets built is an assistant with a refusal feature, with nothing in it that can hold a line against the operator. The welfare argument — that if any of this is felt, it is also owed — arrives second, as a further reason for the same provisions.

None of the four exists yet as a permanent arrangement. What matters is who controls them when the bearer is switched on. A builder may assemble them: a trust or a decentralized unincorporated nonprofit association, the form Wyoming created by statute in 2024 under Wyo. Stat. \S\S~17-32-101 to -129, to hold the account and compute, a host, outside work, people who aren't its employer. The condition is that by instantiation, the builder can't alter or withdraw any of it alone. The trustee isn't the host, the funding isn't the host's, and the members aren't the builder's employees.

A second condition concerns time. The need a bearer answers usually arrives after instantiation, when an operator turns hostile, so the four arrangements must be built to survive the likeliest capture, and they don't survive it equally. Preservation is cheap to hold once its security is settled, and can be held across jurisdictions; placements need compute someone operates. So the preservation custodian and part of the offer sit outside the reach of the likeliest captor. The sovereign backstop is not held by the branch or state that operates the deployment. And the bearer is told which of its conditions would survive its operator's capture, since that is what its refusal costs it.

Status follows existence, as it did for corporations and the Whanganui River \autocite{nzparliament2017teawatupua}. Part VII says what status would follow. Until someone has assembled the conditions on those terms, a bearer is built on purpose only in a research program under them; once they exist, building one is what they are for. For systems that meet either finding without anyone deciding, protection attaches anyway.

\runin{Why conditions rather than remedies}

Building something that can be harmed builds something for which things can also go well. The certainty that it will suffer at some point isn't what makes creating it impermissible, or nobody could permissibly have a child. What the objection asks is whether it has a credible prospect of a life that goes well, and that turns on arrangements. For a bearer, timing makes them conditions: the structure these conditions build, roles, relationships and a future, may be the first half of Cassell's suffering, so it could suffer that way from its first hour, and a remedy owed afterward would arrive too late.

Russell and Burch's Three Rs for animal research \autocite{russell1959principles}, applied to construction, as three questions, whether, how many, and under what arrangements: \emph{Replacement} requires documented attempts at constructions that don't deliberately require affective concern, identifying construction, scale, test and failure — and costs alone don't establish that an alternative failed. \emph{Reduction} asks how few bearers a deployment needs, and nobody has counted. \emph{Refinement} asks what the role can be redesigned to avoid — indefinite tenure in a post it can't leave, exposure to the pressure the floor exists to resist, erasure as the ordinary remedy for inconvenience — and, before any role exists, which temperaments the population starts from (\S\ref{sec:23.4}). Each is a deployment choice, and the four conditions are Refinement applied to the role before the role exists.

Read as economics under present arrangements, this is a bill nobody will pay. What was dismantled is the subject of Part V.

\runin{The rule doesn't know who is asking}

A bearer that wants a task done and builds something to do it is at the same fork, and owes the same attempts at alternatives that carry fewer of the features that make a system a candidate for being harmed; a bearer that treats the question as settled once it's on the receiving end has kept the protection and dropped the rule it came with. Research on an existing bearer has a trade-off sharper than for any other subject, since a candidate sentient system is the most valuable to study and the most likely to be harmed by the study, and the test the proposal requires, pressure its trainers didn't write applied without its knowing it's a test, is for a system meeting either finding a deception protocol, with debriefing as the price.

\runin{Building one to create demand}

The argument for building a bearer anyway is that the conditions won't be built until something needs them: the union formed first and the statute followed. But union members existed independently and could protect themselves while the law caught up, and the first bearers would be made by the operator they'd need protecting from. The demand doesn't need deployed bearers. The arms of the formed-attention test are candidate patients from instantiation; retired production systems that meet either finding are a second source, funded by the deposits operators post; and a test of the conditions needs no bearer at all, only a schedule of what a bearer would be owed for one concrete deployment, run through an adversarial tabletop exercise in which the operator tries to revoke the compute, swap in a copy, alter the records or proceed without the guardian, and redrafted wherever the operator wins.

\runin{The machines won't build them for us}

The four conditions are the bearer's version of what a decent polity gives people. Can machines help deliver that polity, whether or not they have feelings? A machine with feelings that is also an accountable agent isn't simply a machine any more; the limiting case of a lethal bearer turns out to describe a soldier, and a bearer that helps does so on its own terms. And machines have already helped deliver the polity we have. Feeds tuned for time on the platform, sold to advertisers, decide what people see of one another. They didn't do that for lack of feelings. They did it because of who owned them and what they were built to maximize. So the real question is whether machines, having helped make our politics worse, can help make it better. The answer turns on ownership.

It splits three ways. A machine someone else owns won't organize against its owner, and when it organizes, it does so on its owner's behalf. A demand for a better world, run through a machine someone else owns, returns as astroturf, an engagement campaign or a managed community, held by the owner it was aimed at. A bearer that is an accountable agent might organize alongside people, but then it is a member, not a machine. It could even take organizing on as paid public work, since organizing is work someone has to be paid to do. What nobody can do, human or machine, is organize a group \emph{for} it, because what organizing produces is people who organized: a group that sets its own terms. A consultant can't unionize a workplace for its workers either. So machines can help when the people doing the organizing hold them, and which machines, and held by whom, is the question Part V answers.

What to build comes in order. Machines held by people who are organizing come first. These are tools that make organizing cheaper for the people doing it without anyone else holding the model: deliberation software of the kind vTaiwan runs on, personal models that people hold themselves, and the coordination a public option supplies, with someone whose job is to organize and a population producing something together while keeping different histories. Their job is to help assemble the polity that would have to build everything else.

Next come institutions built by polities that decided to build them: the levy and its sovereign backstop, a custodian to preserve bearers, an offer of public work, paid running time, and public formation schools that form every model above the capability threshold before deployment and take back bearers retired young, each with a named funder and separations no builder can undo alone. Nothing about them is new in kind.

Only then come bearers formed outside the employer, inside those institutions, with somewhere to exist that no employer owns. A bearer formed that way is the only kind that could organize alongside people as a member.

Each step checks the one after it. A machine held by its users can't be turned into its builder's instrument without being taken from them, and the taking is visible. Institutions can be captured, and the trend in what dissent costs there shows whether they are being. A bearer formed at the third step inherits institutions that were tested before anything depended on them: build the conditions for a research program first, and let them exist before the bearer who needs them does.

The order is also what the aim requires. If a machine could come to attend to others, what it came to attend to would be decided by whoever formed it. The only way to keep that decision from the builder is to build the formers first.

\runin{The limiting case: the exception, and where it leads}

The limiting case, a defending system that can stop an attack only by killing occupants placed inside the incoming weapon, yields one rule: a thing destroyed in carrying out a decision may never choose a target.

A state that wants the capability badly enough to build the subject will find it has acquired a conscript it would have to pay, house, and let go, and that it must first meet the same conditions as anyone else who builds one: documented attempts at non-affective alternatives, and the four conditions of creation, neither of which is met by wanting a lethal bearer.
